\section{线性变换与基底} \label{sec:vector-base}

在开始本节内容之前, 有一个事情想说一下:
本教程不能安排大段的篇幅来讲解线性代数,
因此这里只用基底的概念对线性变换作一个简略的介绍.
如果读者有余力, 想要更加深入地学习相关内容,
我推荐观看 3B1B 大佬制作的线性代数科普视频:
\hyperlink{https://www.bilibili.com/video/BV1ys411472E/}
{\textcolor{blue}{【官方双语/合集】线性代数的本质 - 系列合集}}.

弹幕制作经常涉及到平移, 旋转, 缩放等几何变换.
空间中的变换将一个位置变换到另一个位置,
我们可以用一个函数来描述它, 该函数的输入和输出都是位置对应的向量.
例如, 平面上的平移变换可以表示为
\[
  f(x,y) = (x+\Delta{x},y+\Delta{y}),
\]
或者用向量运算表示为
\[
  f(\vec r) = \vec r + \Delta{\vec r},
\]
其中$\Delta{\vec r}=(\Delta{x},\Delta{y})$是该平移变换的平移量.

\subsection{线性变换}

线性变换是一类特殊的变换, 它保持向量的线性运算不受影响:

\begin{definition}[线性变换]
  如果空间中的一个变换$f(\vec r)$满足:
  \begin{enumerate}[label=(\alph*)]
    \item 保持加法: 对于空间中任意向量$\vec r_1,\vec r_2$, 有
          \[
            f(\vec r_1+\vec r_2) = f(\vec r_1) + f(\vec r_2).
          \]
    \item 保持数乘: 对于空间中任意向量$\vec r$和任意数值$\lambda$, 有
          \[
            f(\lambda\vec r) = \lambda f(\vec r).
          \]
  \end{enumerate}
  那么我们说这个变换$f$是\emph{线性}的.
\end{definition}

\begin{example}
  放缩变换$f(x,y) = (ax,by)$. 任取实数$\lambda$
  和向量$\vec r_1=(x_1,y_1),\vec r_2=(x_2,y_2)$, 可以验证
  \begin{align*}
    f(\vec r_1+\vec r_2)
     & = f(x_1+x_2,y_1+y_2)          \\
     & = (a\,(x_1+x_2),b\,(y_1+y_2)) \\
     & = (ax_1,by_1)+(ax_2,by_2)     \\
     & = f(\vec r_1)+f(\vec r_2),
  \end{align*}
  \begin{align*}
    f(\lambda\vec r_1)
     & =f(\lambda x_1,\lambda y_1)  \\
     & =(\lambda ax_1,\lambda by_1) \\
     & =\lambda f(\vec r_1).
  \end{align*}
  所以该变换是线性的.
\end{example}

\begin{example}
  平移变换$f(\vec r) = \vec r+\Delta{\vec r}$.
  取两个向量$\vec r_1,\vec r_2$, 计算
  \begin{align*}
    f(\vec r_1 + \vec r_2)
     & = \vec r_1 + \vec r_2 + \Delta{\vec r},  \\
    f(\vec r_1) + f(\vec r_2)
     & = \vec r_1 + \vec r_2 + 2\Delta{\vec r}.
  \end{align*}
  只要平移量$\Delta{\vec r}\neq0$,
  就有$f(\vec r_1 + \vec r_2) \neq f(\vec r_1) + f(\vec r_2)$.
  所以平移变换不是线性的.
\end{example}

\begin{theorem}\label{thm:linear-zero}
  若$f$为线性变换, 则有$f(0) = 0$.
\end{theorem}
\begin{proof}
  任取向量$\vec r$, 有
  $f(0) = f(\vec r) - f(\vec r) = 0$.
\end{proof}

定理\ref{thm:linear-zero}表明,
线性变换必然保持原点位置不变.
常见的变换中, 以原点为中心的旋转变换和放缩变换是线性的,
而平移变换以及不以原点为中心的旋转和放缩是非线性的.

为了进一步讨论线性变换, 我们引入\emph{基底}的概念.

\subsection{空间的基底}

我们说向量在几何上对应只考虑长度和方向不考虑起点位置的有向线段,
在代数上对应由若干个数值构成的数组.
通过直角坐标系上的坐标, 我们把几何上的向量和代数上的向量联系起来.
向量的内核是线性运算, 即加法运算和数乘运算,
我们希望能够从向量线性运算的角度去理解
为什么我们用直角坐标系的坐标去做这样的几何与代数的联系.

回顾一下我们是如何在平面直角坐标系上描述点的位置.
如图\ref{fig:vector-base}左侧,
平面上有一点$P$, 在直角坐标系下的坐标为$(2,3)$
(自然地, 向量$\vv{OP}$在直角坐标系下的坐标也为$(2,3)$).
此时我们作矩形$OAPB$,
$\vv{OP}$可以用向量加法$\vv{OP} = \vv{OA} + \vv{OB}$表示.
同时, 我们取$x$轴正方向的单位向量$\vec i=(1,0)$
以及$y$轴正方向的单位向量$\vec j=(0,1)$,
那么$\vv{OA},\vv{OB}$又可以用向量数乘
$\vv{OA}=2\vec i,\vv{OB}=3\vec j$表示.
于是, 向量$\vv{OP}$可以用$\vec i,\vec j$的线性组合来表示:
\[\vv{OP} = 2\vec i + 3\vec j.\]

所以, 从向量运算的角度来看,
平面直角坐标系的实质就是提供两个(互相垂直的单位)向量$\vec i,\vec j$,
把平面上的所有向量用这两个向量的线性组合表示,
而坐标值$(x,y)$则表示这个线性组合的数乘系数.

\begin{figure}[htbp]
  \centering
  \begin{subfigure}{0.4\textwidth}
    \centering
    \begin{tikzpicture}[]
      \foreach \x in {0,1,2}
      \draw[help lines] (\x,-0.5) -- (\x,3.5);
      \foreach \y in {0,1,2,3}
      \draw[help lines] (-0.5,\y) -- (2.5,\y);

      \draw (0,0) node[below left] {$O$}
      -- (2,0) node[below right] {$A$}
      -- (2,3) node[above right] {$P$}
      -- (0,3) node[above left] {$B$}
      -- cycle;
      \begin{scope}[-latex]
        \draw (0,0) -- (1,0) node[below,background]
        {$\vec i=(1,0)$};
        \draw (0,0) -- (0,1) node[left,background]
        {$\vec j=(0,1)$};
        \draw (0,0) -- (2,3);
      \end{scope}
    \end{tikzpicture}
  \end{subfigure}
  \begin{subfigure}{0.4\textwidth}
    \centering
    \begin{tikzpicture}[x=(30:1),y=(135:1.5)]
      \foreach \x in {0,1,2}
      \draw[help lines] (\x,-0.5) -- (\x,3.5);
      \foreach \y in {0,1,2,3}
      \draw[help lines] (-0.5,\y) -- (2.5,\y);

      \begin{scope}[-latex]
        \draw (0,0) -- (1,0) node[right,background] {$\vec e_1$};
        \draw (0,0) -- (0,1) node[left,background] {$\vec e_2$};
        \draw (0,0) -- (2,3) node[above,background]
        {$\vec r = 2\vec e_1 + 3\vec e_2$};
      \end{scope}
    \end{tikzpicture}
  \end{subfigure}
  \caption{平面的基底}
  \label{fig:vector-base}
\end{figure}

实际上, 我们可以用任意两个不共线的向量
的线性组合来表示平面上的所有向量.
例如, 在图\ref{fig:vector-base}右侧,
$\vec e_1,\vec e_2$是两个不共线的向量,
对于平面上的任一向量(比如图中的$\vec r$)
都可以表示为$\vec e_1,\vec e_2$的线性组合
(比如图中为$\vec r = 2\vec e_1 + 3\vec e_2$).

\begin{definition}[基底]
  设$\vec e_1,\vec e_2,\dots,\vec e_n$是$n$维空间
  中的一组线性无关的向量,
  则对于空间中的任意向量$\vec\alpha$,
  存在唯一一组实数$a_1,a_2,\dots,a_n$使得
  \[
    \vec\alpha = a_1 \vec e_1 + \dots + a_n \vec e_n.
  \]

  此时, 将向量组$(\vec e_1,\dots,\vec e_n)$
  称为空间的一个\emph{基底},
  数组$(a_1,\dots,a_n)$称为向量$\vec\alpha$在该基底下的坐标.
\end{definition}

\begin{definition}[线性无关]
  我们说一组向量$\vec\alpha_1,\dots,\vec\alpha_t$是线性无关的,
  如果这些向量中没有一个向量可以表示为其他向量的线性组合.
  即, \[
    \lambda_1\vec\alpha_1 + \dots + \lambda_n\vec\alpha_n = 0
  \]
  只在$\lambda_1=\dots=\lambda_n=0$时成立.

  $t=2$时, 这意味着$\vec\alpha_1,\vec\alpha_2$不能共线;
  $t=3$时, 这意味着$\vec\alpha_1,\vec\alpha_2,\vec\alpha_3$不能共面.
  依此类推.
\end{definition}

\begin{example}[默认基底]
  为了描述$n$维空间中的位置, 我们通常会有一个默认的直角坐标系,
  每个位置在该坐标系中对应一个坐标.
  我们一般只说这是一个坐标, 而不说这是某个基底下的坐标,
  因为每次都要强调是哪个基底真的很麻烦 (doge).

  这个默认的基底包含$n$个向量,
  每个基向量对应一个坐标轴正方向的单位向量.
  在平面上, 我们用$(\vec i,\vec j)$表示默认基底,
  其中$\vec i=(1,0), \vec j=(0,1)$;
  在三维空间, 我们用$(\vec i,\vec j,\vec k)$表示默认基底,
  其中$\vec i=(1,0,0), \vec j=(0,1,0), \vec k=(0,0,1)$.
\end{example}

\begin{example}[标准正交基]
  \emph{标准正交基}是一类特殊的基底,
  它的各个基向量两两垂直
  (对应名称里的``正交''\footnote{
    垂直是几何上的概念, 表示几何对象之间的直角夹角关系;
    正交是代数上的概念, 表示向量的内积为零.
  }), 且长度都为1 (对应名称里的``标准'').
  例如默认基底就属于标准正交基.

  标准正交基有一些很方便的性质, 举一个例子:
  求一个向量$\vec\alpha$在标准正交基
  $(\vec e_1,\dots,\vec e_n)$下的坐标$(a_1,\dots,a_n)$,
  只需要计算向量内积:
  $a_i = \vec\alpha \cdot \vec e_i$ ($i=1,\dots,n$).
  我们很容易进行证明:
  对于$1\le i,j \le n$, 有
  \[
    \vec e_i \cdot \vec e_j =
    \begin{cases}
      1, & 若\ i = j;    \\
      0, & 若\ i \neq j.
    \end{cases}
  \]
  然后考虑$\vec\alpha = a_1\vec e_1 + \dots + a_n\vec e_n$, 可得
  \begin{align*}
    \vec\alpha \cdot \vec e_i
     & = a_1\vec e_1\cdot\vec e_i + \dots
    + a_n\vec e_n\cdot\vec e_i            \\
     & = a_i\vec e_i\cdot\vec e_i         \\
     & = a_i.
  \end{align*}
\end{example}

\subsection{用基底表示线性变换}

回顾线性变换的性质, 由于线性变换保持加法和数乘,
一旦我们知道了线性变换$f$对默认基底的作用,
我们就能完全确定该变换对空间中所有向量的作用.

假设$(\vec e_1,\dots,\vec e_n)$是空间的一个基底,
$\vec a = a_1\vec e_1 + \dots + a_n\vec e_n$
是空间中的任意向量. 根据线性性质:
\begin{align*}
  f(\vec a) & = a_1f(\vec e_1) + \dots + a_nf(\vec e_n).
\end{align*}

我们发现, 当$\vec a$在默认基底
$(\vec\varepsilon_1,\dots,\vec\varepsilon_n)$
下的坐标为$(a_1,\dots,a_n)$时,
变换后的$f(\vec a)$在基底
$(f(\vec\varepsilon_1),\dots,f(\vec\varepsilon_n))$
下的坐标也为$(a_1,\dots,a_n)$.
\footnote{
  严格来说$(f(\vec\varepsilon_1),\dots,f(\vec\varepsilon_n))$
  不一定是基底, 因为我们不能保证它们线性无关.
}
于是我们可以用向量组
$(f(\vec\varepsilon_1),\dots,f(\vec\varepsilon_n))$
表示线性变换$f$, 记作
\[
  f = (\vec f_1,\dots,\vec f_n),
\]
其中$\vec f_i = f(\vec\varepsilon_i)$ ($i=1,\dots,n$).
此时对于空间中任一向量$(a_1,\dots,a_n)$, 有
\[
  f(a_1,\dots,a_n) = a_1 \vec f_1 + \dots + a_n \vec f_n.
\]

\subsection*{习题}

\begin{enumerate}
  \item \comment{旋转公式v2}
        设平面上有一点$(x,y)$, 将该点绕原点旋转$\phi$角
        (以逆时针为正方向), 得到的点的坐标$(x',y')$怎么表示?
        这就是二维旋转公式要解决的问题.

        上述的旋转变换是一个线性变换.
        请你用向量组表示该变换, 并由此证明二维旋转公式:
        \begin{align*}
          x' & = x\cos\phi + y\cos(\phi+90\degree), \\
          y' & = x\sin\phi + y\sin(\phi+90\degree).
        \end{align*}
  \item \comment{坐标系变换v2**}
        我们可以用平面基底建立一种新的坐标系.
        类似于直角坐标系, 指定平面上的一点$A$为原点,
        以及两个不共线的向量$\vec e_1,\vec e_2$.
        对平面上任意一点$P$, 存在唯一一组实数$a,b$使得
        \[
          \vv{AP} = u\vec e_1 + v\vec e_2.
        \]
        这类由原点和基底向量定义的坐标系称为仿射坐标系,
        记作$(A;\vec e_1,\vec e_2)$,
        $(u,v)$称为点$P$在该坐标系下的仿射坐标.

        本题探讨默认的直角坐标系$(O;\vec i,\vec j)$
        与某个仿射坐标系$(A;\vec e_1,\vec e_2)$之间的坐标转换.
        已知该仿射坐标系的原点$A$和基底$\vec e_1,\vec e_2$
        在默认直角坐标系下的坐标.
        设平面上一点$P$的直角坐标为$\vec r=(x,y)$,
        在$(A;\vec e_1,\vec e_2)$下的仿射坐标为$\vec s=(u,v)$.

        由点$P$的仿射坐标$\vec s=(u,v)$
        计算直角坐标$\vec r=(x,y)$是很容易的.
        根据仿射坐标系的定义,
        \[
          \vec r = A + \vv{AP} = A + u\vec e_1 + v\vec e_2,
        \]
        代入$A,\vec e_1,\vec e_2$的直角坐标,
        即可计算出直角坐标$\vec r$.

        由直角坐标$\vec r=(x,y)$
        计算仿射坐标$\vec s=(u,v)$则比较麻烦,
        我们从线性变换的角度着手.

        设线性变换$f = (\vec e_1,\vec e_2) = ((a,b),(c,d))$,
        则$\vv{AP} = f(\vec s)$.
        只要能求出$f$的逆变换$f^{-1}$, 就可以计算出仿射坐标
        $\vec s = f^{-1}(\vv{AP}) = f^{-1}(\vec r - A)$.

        实际上, $f^{-1}$也是一个线性变换, 并且对于任意向量$\vec\alpha$,
        $f^{-1}$总是把$f(\vec\alpha)$变换到$\vec\alpha$.

        问题: 已知$f=(\vec e_1,\vec e_2)=((a,b),(c,d))$,
        请计算$f^{-1}$.
\end{enumerate}
